% Source-derived chapter 04: A bound for Betti numbers
% Original source: bounds-stalks-perverse-sheaves-characteristic-p-shende
\chapter{A bound for Betti numbers}
\label{bounds-stalks-perverse-sheaves-characteristic-p-shende-chapter-04}
\source{bounds-stalks-perverse-sheaves-characteristic-p-shende}

\begin{remark}[Source section]
\label{bounds-stalks-perverse-sheaves-characteristic-p-shende-chapter-04-source}
\source{bounds-stalks-perverse-sheaves-characteristic-p-shende}
\source{bounds-stalks-perverse-sheaves-characteristic-p-shende}
This chapter reproduces the corresponding section of the retrieved
LaTeX source, including its definitions, statements, examples, and
proofs. It has not yet been translated into Lean.
\notready
\end{remark}

% The source section title was: A bound for Betti numbers
\begin{lemma}
\source{bounds-stalks-perverse-sheaves-characteristic-p-shende}
\notready\label{transversality-check} Let $f: X \to Y$ be a smooth morphism of smooth varieties with $X$ of dimension $n$ and $Y$ of dimension $n-m$. Let $C$ be a closed conical subset of the cotangent bundle $T^* X$ of $X$ with all irreducible components of dimension  $n$. Let $y$ be a point in $Y$ and let $i$ be the inclusion of $f^{-1}(y)$ into $X$, so that we have a Cartesian square. \begin{equation}\label{i-diagram}\begin{tikzcd} f^{-1}(y) \arrow[r,"i"] \arrow[d] & X \arrow[d,"f"] \\ y \arrow[r] & Y \end{tikzcd}\end{equation}
\source{bounds-stalks-perverse-sheaves-characteristic-p-shende}

 If $f$ is $C$-transversal and the fibers of the composition $C \to X \to Y$ have dimension $m$, then $i$ is properly $C$-transversal. \end{lemma}

\begin{proof} Because $f$ is $C$-transversal, the inverse image of $C$ by $df: X \times_Y T^* Y \to T^* X$ is a subset of the zero-section. Hence the intersection of $C$ with the image of $df$ is a subset of the zero-section, as only nonzero points are sent to nonzero points by $df$. The image of $df$ in $T^* X$ consists of $1$-forms that are pulled back from $Y$, i.e. one-forms that are transverse to the fibers of $X$, which are exactly those one-forms in the kernel of $di: f^{-1}(y) \times_X T^* X \to T^* f^{-1}(y)$. Hence $i$ is $C$-transversal.

$i^* C = f^{-1}(y) \times_X C =  y \times_Y C$ is exactly a fiber of the composition $C \to Y$, and thus the claim that it has dimension $\dim X - \dim Y = \dim (f^{-1}(y))$ verifies that $i$ is properly $C$-transversal.

\end{proof}


\begin{lemma}
\source{bounds-stalks-perverse-sheaves-characteristic-p-shende}
\notready\label{nearby-formula} Let $X$ be a smooth variety and $Y$ a smooth curve, both over a perfect field $k$. Let $K$ be an object in $D^b_c(X,\mathbb F_\ell)$.
\source{bounds-stalks-perverse-sheaves-characteristic-p-shende}

Let $CC'(K)$ be $CC(K)$ with any occurrence of the cotangent space at $x$ removed, and let $SS'(K)$ be $SS(K)$ with any occurrence of the cotangent space at $x$ removed. 

Let  $f: X \to Y$ be a smooth projective morphism that is $SS'(K)$-transversal and such that the fibers of $SS'(K)$ over $Y$ have dimension $\dim X-1$ . Let $y=f(x)$, let $Z =f^{-1}(y)$ and let $i$ be the inclusion of $Z$ into $X$, as in Diagram \ref{i-diagram}.  Then 
\[i^! CC'(K) = CC ( R \Psi_{f} K) \]
where we view the nearby cycles relative to $f$ as a complex of sheaves on $f^{-1}(Y)$. \end{lemma}

Note that $SS(K)$ here is a union of irreducible varieties of dimension $\dim X$ and $CC(K)$ is a $\mathbb Z$-linear combination of irreducible varieties of dimension $\dim X$. When we refer to removing the cotangent space at $X$, an irreducible variety of dimension $\dim X$, we mean removing this term from the $\mathbb Z$-linear combination or the union, if it appears, but leaving all other terms.


\begin{proof} Because $f$ is $SS'(K)$-transversal, it is $SS(K)$-transversal away from $x$, so $K$ is locally acyclic away from $f$ by the definition of the singular support, and thus $R \Phi_f K$ vanishes away from $x$, so $R \Psi_f K = i^* K$ away from $x$.

Furthermore, $i$ is properly $SS'(K)$-transversal by Lemma \ref{nearby-formula}.

Then by \cite[Theorem 7.6]{saito1} \[  CC ( R \Psi_{f} K) = CC( i^* K ) = i^! CC(K) = i^! CC'(K) \] away from $x$. 

Hence \[i^! CC'(K) -  CC ( R \Psi_{f} K) \] is a cycle on the cotangent bundle of $Z$ supported inside the cotangent space at $x$. Because these cycles are rational linear combinations of irreducible closed sets of dimension $\dim Z$, the difference is a multiple of the cotangent space at $x$. Because the cotangent space at $x$ has nonzero intersection number with the zero-section $Z$ of $T^* Z$, to check that \[ i^! CC'(K) = CC(R\Psi_f K),\] it suffices to check \[ (i^! CC'(K) , Z)_{T^*Z}=   (CC ( R \Psi_{f} K) , Z)_{T^*Z} . \] By the index formula \cite[Theorem 7 .13]{saito1},  \[ (CC ( R \Psi_{f} K) , Z)_{T^*Z} =  \chi(Z, R \Psi_f K) =\chi ( f^{-1}(\eta),K)\] for $\eta$ the generic point of $Y$. 

By definition, \[i^! CC'(K) = - (di)_* i^*  CC'(K).\] We have \cite[Proposition 8.1.1(c)]{Fulton}
\[  ((di)_* i^* CC'(K), Z)_{T^*Z} = (i^* CC'(K),  (di)^* Z)_{T^*X \times_X Z} = (CC'(K), i_* (di)^* Z) _{T^*X} .\]
For the first identity, this uses the fact that $di$ is finite on the support of $i^* CC'(K)$ and for the second identity this uses the fact that $i$ is a closed immersion, hence finite.

Now $(di)^ * Z \subseteq  T^* X \times_X Z$ consists of one-forms transverse to $Z$, so $i_* (di)^* Z$ is the conormal bundle of $Z$ inside $X$, $N^*Z$. 

As a point $y' \in Y$ varies, the conormal bundle to $f^{-1}(y')$ varies in a smooth family. To check that the intersection number  \[ (CC'(K), N^*f^{-1} (y') )_{T^*X} \] is constant, it suffices to check that the the intersection $ CC'(K) \cap N^*f^{-1} (y') $, viewed as a family of closed subsets parameterized by $y' \in Y$, and hence viewed as a scheme mapping to $Y$ given the induced reduced subscheme structure, is proper over $Y$. This is true because, as $f$ is $SS'(K)$-transversal, this intersection is contained in the zero-section, hence is proper.

The same is true for any other $y'$, and these conormal bundles vary in a smooth family, so this intersection number for $y$ is equal to the intersection number for any $y'$, and in particular for the generic point $\eta$. Let $i_\eta$ be the inclusion of the generic fiber of $f$ into $X$, then \[( i^! CC'(K), Z)_{T^*Z}  = - (CC'(K),  N^* f^{-1}(y))_{T^*X} = (CC'(K), N^* f^{-1}(\eta))_{T^*X}\] 
\[ = (i_\eta^! CC'(K), f^{-1}(\eta))_{T^* f^{-1}(\eta)}  =  (i_\eta^! CC(K), f^{-1}(\eta))_{T^* f^{-1}(\eta)} =  ( CC(i_\eta^* K) , f^{-1}(\eta))_{T^* f^{-1}(\eta)}\] \[= \chi(f^{-1}(\eta),K) \] as desired, where the first equality summarizes the previous calculations.\end{proof}

\begin{lemma}
\source{bounds-stalks-perverse-sheaves-characteristic-p-shende}
\notready\label{genericity-lemma} Let $X$ be a smooth variety embedded in projective space $\mathbb P^n$. Let $C$ be a closed conical subset of the cotangent space of $X$ of dimension $\dim X$. Let $x \in X$ be a point such that $C$ does not contain the cotangent space of $x$.
\source{bounds-stalks-perverse-sheaves-characteristic-p-shende}

Let $\overline{X} \subseteq X \times \mathbb P^1$ be a general pencil of conic sections of $X$, parameterized by $\mathbb P^1$.  Let $p$ and $q$ be the natural projections in the below commutative diagram.\begin{equation}\label{pq-diagram}\begin{tikzcd} \overline{X}  \arrow[ddr,"p"'] \arrow[drr,"q"] \arrow[dr]  \\  &X \times \mathbb P^1 \arrow[r] \arrow[d] & \mathbb P^1 \\ & X \end{tikzcd} \end{equation}
\source{bounds-stalks-perverse-sheaves-characteristic-p-shende}

Then $p$ is properly $C$-transversal, $q$ is $p^{\circ} C$-transversal in a neighborhood of the unique conic in the pencil containing $x$, and the fibers of $p^\circ C$ over $\mathbb P^1$ are $\dim X-1$-dimensional in a neighborhood of the conic containing $x$. \end{lemma}

\begin{proof}

Let $Y$ be the base locus of this pencil of conics, which since the pencil is generic, is a smooth subscheme of codimension $2$ (by Bertini's theorem). We can view $p$ as the blow-up of $X$ along $Y$.

First we check that $p$ is $C$-transversal. The map $p$  is \'{e}tale, and automatically $C$-transversal, away from $Y$, and at each point over $Y$, $dp^{-1} ( \{0\})$ is one-dimensional and contained in the conormal space of $Y$. Thus to check that $p$ is $C$-transversal, it suffices to check that no point in $C$ consists of a point in $Y$ and a nonzero vector transverse to $Y$.  For each pair of a point and nonzero cotangent vector in $C$, the condition that the point be contained in $Y$ is a codimension $2$ condition on the pencil of conics, and the condition that the vector be transverse to $Y$ is a codimension $\dim X-2$ condition on the pencil of conics. Because the space of pairs of a point and a nonzero cotangent vector contained in $C$, up to dilation of the cotangent vector, is $\dim X-1$-dimensional, this occurring for any point is a codimension $1$ condition, hence is not generic. 


Next we check that $p$ is properly $C$-transversal.  Because $\dim \overline{X} =\dim X$ and the fibers of $p$ have dimension at most one, $\dim p^* C = \dim C = \dim X = \dim \overline{X}$ unless the base of some irreducible component of $C$ lies entirely in $Y$. For any given variety, a generic $Y$ does not contain it, so this does not happen, and $p$ is properly $C$-transversal.

 Let $t$ in $\mathbb P^1$ be such that $x \in \overline{X}_t$. Now we check that $q$ is $p^\circ C$-transversal in a neighborhood of $X_t$.  Because $dq^{-1}  (p^\circ C)$ is a closed conical subset of $ \overline{X} \times_{\mathbb P^1} T^* \mathbb P^1$, and a closed conical subset being contained in the zero section is an open condition, $q$ being $p^\circ C$-transversal is an open condition.
 
 Thus it suffices to check that the restriction of $dq^{-1} (p^\circ C)$ to $X_t$ is contained in the zero section. Equivalently, we fix a nonzero one-form $\omega_0$ on $\mathbb P^1$ at $t$, so that $dq(\omega_0)$ generates the one-dimensional image of $dq$, and check that $dq(\omega_0)_z \notin p^\circ C$ at each point $y\in X_t$. 
 
%The map $q$ is $p^\circ C$-transversal at a point $z \in X_t$, unless  $dq^{-1} ( ( p^\circ C)_{z})$ contains a nonzero vector. Because $q$ is a map to a one-dimensional variety, the image of $dq$ is one-dimensional, generated by a single vector $dq(\omega)$, and $q$ is $p^\circ C$-transversal unless $dq(\omega) \in p^\circ C$.

If $y\in Y$, the image of $dp$ and $dq$ intersect only at zero, so $dq(\omega_0)_z \notin p^\circ C$. If $z\notin Y$,  $dq(\omega_0 )$ is the conormal vector to  $X_t$, because $X_t$ is a level set of $q$. Thus, to check that  $q$ is generically $p^\circ C$-transversal in a neighborhood of $\overline{X}_t$, it suffices to check that, for a general conic $\overline{X}_t$ through $x$, the conormal bundle to $\overline{X}_t$ never contains a pair of a point and a nonzero cotangent vector in $C$. 

For each point $y$ and nonzero cotangent vector $\omega_1$ in $C$, with $y\neq x$, the conics through $x$ whose conormal bundles contain $(y, \omega_1)$ form a codimension $\dim X$ subset of the conics through $x$, because this is a codimension one condition on the value of the conic at $y$ and a codimension $\dim X-1$ condition on the derivative of the conic at $y$, and the derivatives at $y$ are independent of the condition that the conic pass through $x$. Because the space of pairs of a point $y$ and a nonzero cotangent vector $\omega_1$ contained in $C$, up to dilation of $\omega_1$, is $\dim X-1$-dimensional, this is a codimension $1$ condition and is not generic. Over the point $x$, the conormal bundle of.a generic conic is a general cotangent line, so it remains to check that $C$ does not contain a general point of the cotangent space at $x$, which holds because we have assumed that the cotangent space at $x$ does not lie in $C$.

Because each irreducible component of $p^\circ C$ has dimension $\dim X$, the fibers over $\mathbb P^1$ have dimension $\dim X -1$ unless some irreducible component is contained entirely in one fiber, i.e. in a single conic in the pencil.  For a generic pencil of conics, the only variety that is necessarily contained in one fiber of the pencil is a single point, and because $C$ does not contain the cotangent space of $x$, none of these points are $x$, and so they will not generically be in the same fiber as $x$, and thus we can remove the fibers containing these points from our chosen neighborhood.
\end{proof}

\begin{lemma}
\source{bounds-stalks-perverse-sheaves-characteristic-p-shende}
\notready\label{final-nearby-formula} Let $X$ be a smooth variety embedded in projective space $\mathbb P^n$ over a perfect field $k$. Let $K$ be an object of $D^b_c(X, \mathbb F_\ell)$.
\source{bounds-stalks-perverse-sheaves-characteristic-p-shende}

Let $CC'(K)$ be $CC(K)$ with any occurrence of the cotangent space at $x$ removed, and let $SS'(K)$ be $SS(K)$ with any occurrence of the cotangent space at $x$ removed. 


Let $\overline{X} \subseteq X \times \mathbb P^1$ be a general pencil of conic sections sections of $X$, parameterized by $\mathbb P^1$.  Let $p$ and $q$ be the projections $\overline{X} \to X$ and $\overline{X} \to \mathbb P^1$, respectively, as in Diagram \ref{pq-diagram}. 

Let $\overline{x} = p^{-1}(x)$, which, because the pencil is generic, is a single point, and let $y= q(\overline{x})$. Let $i : q^{-1}(y) \to \overline{X}$ be the inclusion of the fiber over $y$ into $\overline{X}$, so that we have the commutative diagram. \begin{equation}\label{pqi-diagram}\begin{tikzcd} \overline{X}  \arrow[ddr,"p"'] \arrow[drr,"q"] \arrow[dr]  &  &  q^{-1}(y) \arrow[ll, "i"'] \arrow[dr] \\  &X \times \mathbb P^1 \arrow[r] \arrow[d] & \mathbb P^1 & y \arrow[l]  \\ & X \end{tikzcd} \end{equation}
\source{bounds-stalks-perverse-sheaves-characteristic-p-shende}

Then

\begin{enumerate}

\item \[ CC ( R \Psi_{q} p^*K) = i^! p^! CC'(K)\]

\item  $R\Phi_q p^* K$ is supported at $x$.

\item  \[-\operatorname{dimtot}  \left( R\Phi_q ( p^* K) \right)_x  \] is the multiplicity of the cotangent space at $x$ in $CC(K)$.

\item If $K$ is perverse, then  $\left( R\Phi_q ( p^* K) \right)_x $ is supported in degree $-1$.

\end{enumerate}
  \end{lemma}

\begin{proof} To obtain (1), we apply Lemma \ref{nearby-formula} to $p^* K$ and $q$. By Lemma \ref{genericity-lemma}, $p$ is properly $SS'(K)$-transversal, and it is \'{e}tale at $x$ so it is properly $SS(K)$-transversal, so by \cite[Theorem 6.6]{saito1}, \[CC(p^* K) = p^!(CC(K))\] and thus \[CC'(p^* K) = p^!(CC'(K)).\] Then by Lemma \ref{genericity-lemma}, $p^* K$ satisifies the conditions of Lemma \ref{nearby-formula}, so \[ CC ( R \Psi_{q} p^*K) = i^! CC'(p^* K) = i^! p^! CC'(K) \] as desired. 

To obtain (2), by Lemma \ref{genericity-lemma}, $q $ is $ p^\circ SS(K)= SS(p^* K)$-transversal in a neighborhood of $x$, minus $x$, hence  $p^* K$ is locally $q$-acyclic away from $x$  by the definition of the singular support,  and thus $R \Phi_q p^* K$ is supported at $x$.



For (3), to calculate $R \Phi_q p^*K $, we use again the fact that $q $ is $ SS'(p^* K)$-transversal, so it is $SS(p^*K)$-transversal away from $x$, and thus $x$ is at most an isolated characteristic point, so by the definition of the characteristic cycle \[ -\operatorname{dimtot} (R \Phi_q p^* K)_x = (CC (p^* K), (dq)^* \omega)_x\] where $\omega$ is nonvanishing one-form on an open neighborhood of $y$ in $\mathbb P^1$.

Because $q$ is $SS'(p^*K)$-transversal, the only irreducible component of $SS(p^*K)$ which intersects $(dq)^* \omega $ is the cotangent space $N^* X$ at $x$. Because $(dq)^* \omega$ is a section of the cotangent bundle, $(N^*x, (dq)^* \omega)= 1$, so $(SS(p^*K), (dq)^* \omega)$ is the multiplicity of $N^*x$ in $CC(p^K)$, which is also the multiplicity of $N^* x$ in $CC(K)$.
% the cotangent space at any point is $1$, so its intersection number with $CC(p^* K)$ is the multiplicity of the cotangent space in $CC(p^* K)$, which is also its multiplicity in $CC(K)$.

For (4), because $p^*K$ is perverse near $x$, $p^* K[-1]$ is perverse near $x$ when restricted to the generic fiber of $q$, and so by the theorem of Gabber \cite[Corollary 4.6]{Illusie}, $R\Phi_q ( p^* K)[-1]$ is perverse (near $x$, and thus everywhere, because it vanishes elsewhere). Because it is perverse and supported at a single point, it is supported in degree $0$, and then the unshifted version is supported in degree $[-1]$.

\end{proof}




\begin{lemma}
\source{bounds-stalks-perverse-sheaves-characteristic-p-shende}
\notready\label{multiplicity-hypersurface}  Let $X$ be a smooth variety embedded in projective space $\mathbb P^n$ over a perfect field $k$. Let $x$ be a point of $X$. Let $C$ be a conical cycle in the cotangent bundle of $X$.  Let $C'$ be $C$ minus any occurrence of the cotangent space at $x$.  Let $\inhyper{X}$ be the intersection of $X$ with a generic conic through $x$, let $j: \tilde{X} \to X$ be the inclusion, and let $\tilde{C} =-  j^{!} C'$.  
\source{bounds-stalks-perverse-sheaves-characteristic-p-shende}


Then for $i>0$, the $i$th polar multiplicity of $C$ at $x$ equals the $i-1$st polar multiplicity of $\inhyper{C}$ at $x$, and for $i=0$, the $i$th polar multiplicity of $C$ at $X$ equals the multiplicity of the cotangent space at $x$ in $C$. \end{lemma}


\begin{proof} We split into three cases: $i=0$, $0< i < \dim X$, $i=\dim X$.

For $i=0$, in the definition of polar multiplicity we can let $Y =X$, with $V$ a rank one subbundle of the cotangent bundle, so $\mathbb P(V)$ is simply a section of $\mathbb P(T^* X)$. If we choose a general section, then the only irreducible component of $\mathbb P(C)$ it intersects at $x$ is the fiber over $x$, which it intersects with multiplicity the multiplicity of that fiber, which is the multiplicity of the cotangent space at $x$ in $C$.

For $0< i < \dim X$, let $Y$ be a general smooth subvariety of $\inhyper{X}$ of codimension $i-1$ passing through $x$ and let $V$ be a general $i$-dimensional sub-bundle of $T^* \inhyper{X}$ over $Y$. Let $V$ be the inverse image of $\inhyper{V}$ in the cotangent bundle of $X$. By Definition \ref{polar-multiplicity-2}, it suffices to check that \[( \mathbb P( \inhyper{C}), \mathbb P(\inhyper{V}))_{\mathbb P( T^* \inhyper{X}),x}  =(  \mathbb P(V) , \mathbb P(C) )_{\mathbb P(T^* X),x}\] and that $Y,V$ satisfies the conditions of Lemma \ref{intersection-comparison} for $X,C$.

Let $s \colon \inhyper{X} \to  \mathbb P(T^* X \times_X \inhyper{X})$ be the section of $\mathbb P(T^* X \times_X \inhyper{X})$ corresponding to the conormal line of $\inhyper{X}$. After replacing $X$ by a suitable neighborhood of $x$, we have a commutative diagram.
\begin{equation}  \begin{tikzcd}  \mathbb P( C' ) \times_X \tilde{X}  \arrow[d,"z" ]  \arrow[r, "t" ] & \mathbb P ( T^* \tilde{X})  \arrow[r,"e"] & \tilde{X} \\
  \mathbb P(T^*X \times_X \inhyper{X}  ) - s(\inhyper{X}) \arrow[r,"u" ] & \operatorname{Bl}_{ s( \inhyper{X} ) } \mathbb P(T^*X \times_X \inhyper{X}  ) \arrow[u,"\rho"] \arrow[r, "b"] &  \mathbb P(T^*X \times_X \inhyper{X}  )  \arrow[u] & \tilde{X} \arrow[l,"s" ] \arrow[ul] \end{tikzcd} \end{equation}
In this diagram, the maps $u, \rho, b$ arise from projective geometry - the blow up of a projective bundle at a section $s$ admits a projection map $\rho$ to the projectivization of the quotient vector bundle, and a map $u$ from the open complement of this section. The inclusion $z\colon   \mathbb P( C' ) \times_X \tilde{X}  \to \mathbb P(T^*X \times_X \inhyper{X}  ) - s(\inhyper{X}) $ exists because $s(x)$ is a general point of $\mathbb P( (T^*X)_x)$ and $C'$ does not contain the whole cotangent space of $X$ at $x$, so $s(x)$ is not contained in $\mathbb P(C')$ and so the image of $s$ is disjoint from $\mathbb P(C')$ over a neighborhood of $X$. We define $t$ as $\rho\circ u \circ x$ to make the diagram commute. 

Because $e\circ t$ is proper and $e$ is separated, $t$ is proper. 
 
By definition, $\tilde{C}$ is the pushforward from $T^* X \times_X \tilde{X}$ to $T^* \tilde{X}$ of the restriction of $C$ from $T^* X$ to $T^* X \times_X \tilde{X}$. Because this pushforward and pullback are compatible with taking quotients by $\mathbb G_m$, we have \[  \mathbb P( \inhyper{C})  = (\rho \circ u) _*  ( \mathbb P(C') \times_X \inhyper{X}) \] and \[(\rho\circ u)^* \mathbb P(\inhyper{V}) =(b \circ u)^*  \mathbb P(V)\]
so \cite[Proposition 8.1.1(c)]{Fulton} \[ ( \mathbb P( \inhyper{C}), \mathbb P(\inhyper{V}))_{\mathbb P(T^* \inhyper{X}), x}   
=  (  \mathbb P(C') \times_X \inhyper{X}  , \mathbb P(V))_{\mathbb P(T^* X)\times_{X} \inhyper{X},x }  = (\mathbb P(C'), \mathbb P(V))_{\mathbb P(T^* X)} \] where the pullback along the open immersion $(b\circ u)$ does not affect the intersection number since the intersection locus is a closed subset of $  \mathbb P(T^*X \times_X \inhyper{X}  ) - s(\inhyper{X})$. 




%There is a projection map \[\rho: \mathbb P(T^*X \times_X \inhyper{X}  ) - s(T^*\inhyper{X})  \to \mathbb P(T^* \inhyper{X}) \] where $s$ is the section of $\mathbb P(T^* X \times_X \inhyper{X})$ corresponding to the conormal line of $\inhyper{X}$.  
%
%Because $s(x)$ is a general point of $\mathbb P( (T^* X)_x)$, and $C'$ does not contain the whole cotangent space of $X$ at $x$,  $s(x)$ is not contained in $\mathbb P(C')$, and so the image of $s$ is disjoint from $\mathbb P(C')$ over a neighborhood of $x$. Hence $\rho$ is well-defined on $C' \times_X \inhyper{X}$. In addition, we can see that $\rho$ is proper when restricted to $\mathbb P(C' )\times_X \inhyper{X}$, because it extends to a proper map from the blow-up of $\mathbb P(T^*X \times_X \inhyper{X}  ) $ at $ s(T^*\inhyper{X})$ and $\mathbb P (C') \times_X \inhyper{X}$ remains closed inside that blow-up.
%
%By definition, $\tilde{C}$ is the pushforward from $T^* X \times_X \tilde{X}$ to $T^* \tilde{X}$ of the restriction of $C$ from $T^* X$ to $T^* X \times_X \tilde{X}$. Because this pushforward and pullback are compatible with taking quotients by $\mathbb G_m$, we have \[  \mathbb P( \inhyper{C})  = \rho_*  ( \mathbb P(C') \times_X \inhyper{X}) \] and \[\rho^* \mathbb P(\inhyper{V}) = \mathbb P(V)\]
%
%
%
%so \cite[Proposition 8.1.1(c)]{Fulton} \[ ( \mathbb P( \inhyper{C}), \mathbb P(\inhyper{V}))_{\mathbb P(T^* \inhyper{X}), x}   
%=  (  \mathbb P(C') \times_X \inhyper{X}  , \mathbb P(V))_{\mathbb P(T^* X)\times_{X} \inhyper{X},x }  = (\mathbb P(C'), \mathbb P(V))_{\mathbb P(T^* X)} \]


Finally we have \[ (\mathbb P(C'), \mathbb P(V))_{\mathbb P(T^* X) } =(\mathbb P(C), \mathbb P(V))_{\mathbb P(T^* X) }\] because the difference between $\mathbb P(C)$ and $\mathbb P(C')$ is the fiber over $x$, and because $Y$ has codimension at least one we can perturb $\mathbb P(V)$ to not intersect this fiber.

Next we check the strict transform condition. The exceptional fiber of the blowup of $ T^* X$ at the cotangent space at $x$ is $\mathbb P ((TX)_x ) \times \mathbb P ( (T^* X)_x)$.  Let $\omega \in (T^*X)_x$ be a conomal vector to $X^*$. Note that $\omega$ is a general vector in $(T^*X)_X$, $(TY)_x$ is general among all $\dim X-i$-dimensional subspaces of $(TX)_x$ perpendicular to $\omega$, and $V_x$ is general among $i+1$-dimensional vector subspaces of $(T^*X)_x$ containing $\omega$. Let $Z$ be the intersection of the strict transform of $\mathbb P(C)$ with the exceptional divisor. We must check that the intersection of $ \mathbb P (  (TY)_x) \times P (V_x)$ with $Z$ vanishes.

The dimension of $Z$ is $\dim X-2$ because $Z$ has dimension one less than $\mathbb P(C)$, which itself has dimension one less than $C$, which has dimension $\dim X$.
%Let $\mathcal M_3$ be the moduli space of triples $  ( \omega, (TY)_x, V_x)$ with $\omega \in (T ^*X)_x$ a nonzero vector, $(TY)_x$ a $\dim X-i$-dimensional subspace of $(TX)_x$ perpendicular to $\omega$, and $V_x$ is a $i+1$-dimensional vector subspaces of $(T^*X)_x$ containing $\omega$. Let $\mathcal M_2$ be the moduli space of pairs $v_1 \in P ( (T^*X)_x) , v_2 \in \mathbb P( (TX)_x)$. Let $\mathcal M_5 \subset \mathcal M_3 \times \mathcal M_2$ be the moduli space satisfying the additional conditions $v_1 \in \mathbb P ( (TY)_x), v_2 \in \mathbb P(V_x)$.  Let $\pi_3: \mathcal M_5 \to \mathcal M_3$ and $\pi_2: \mathcal M_3 \to \mathcal M_2$ be the projections
%Because $\dim Z= \dim X -2$, it suffices to check that the fiber in $\mathcal M_5$ over any point in $\mathcal M_3$ has codimension at least $\dim X -1 $ in $\mathcal M_2$. 
So it suffices to prove that for $(v_1,v_2) \in Z \subset \mathbb P ( (T^*X)_x) \times \mathbb P( (TX)_x)$, the codimension of the space of triples $( \omega, (TY)_x, V_x)$ such that $v_1\in (TY)_x , ,v_2 \in V_x $ inside the space of all triples $(\omega, (TY)_x, V_x)$ is at least $\dim X-1$. 



%This codimension is constant on orbits of the  $GL_{\dim X}$ action on pairs $v_1,v_2$. Because there are only two orbits, the codimension depends only on whether $v_1 \cdot v_2 =0$.  On the open orbit where $v_1 \cdot v_2 \neq 0$, the codimension is the dimension of the moduli space of tuples $v_1, v_2, \omega, (TY)_x, V_x$ 

%\[ \dim \mathbb P ( (TX)_x) - \dim \mathbb P ( (TY)_x) + \dim \mathbb P ((T^*X)_x)  - \dim (\mathbb P ( V_x)) =  \dim X - (\dim X-i)  + \dim X- (i+1) = \dim X- 1 .\] 



The pairs  $(\omega,  V_x)$ such that $\omega, v_2 \in V_x$ have codimension $\dim X - (i+1)$ in the space of pairs $(\omega, V_x)$ with $ \omega \in V_x$. (We can ignore $\omega$ for this calculation).

To have $v_1 \in (TY)_x$ and $(TY)_x$ perpendicular to $\omega$, we must have $v_1 \cdot \omega =0$. This is a codimension $1$ condition on $(\omega, V_x)$. To check this, note that the space of pairs $\omega, V_x$ with  $\omega, v_2 \in V_x$ is irreducible, so it suffices to check the function $v_1 \cdot \omega$ is not identically zero. We can do this by choosing $\omega$ with $v_1 \cdot \omega$ nonzero and then choosing $V_x$ to contain $\omega$ and $v_2$, using $\dim V_x= i+1 \geq 2$ since we have assumed $i>0$. 


Over each pair $(\omega, V_x)$ with  $\omega, v_2 \in V_x$ and $v_1 \cdot \omega = 0$,  the $\dim X-i$-dimensional spaces $(TY_x)$ that contain $v_1$ and are perpendicular to $\omega$ have codimension $i-1$ among all $\dim X-i$-dimensional spaces perpendicular to $\omega$, for a total codimension of $\dim X - 1$.

This completes the case $0 < i < \dim X$.


 


%(The intersection of the exceptional fiber with the strict transform of $\mathbb P(V)$ is $\mathbb P ( (TY)_x ) \times \mathbb P ( V_x)$. Note that $(TY)_x$ is a general $\dim X-i$-dimensional vector subspace of $(TX)_x$ contained in the tangent space of $\inhyper{X}$, and $V_x$ is a general $i+1$-dimensional vector subspace of $(T^* X)_x$ containing the conormal line of $\inhyper{X}$.
%
%So it suffices to prove that, for a general vector $\omega$ in $(T^*X)_x$, $ (TY)_x$ a general $\dim X-i$-dimensional subspace of $(TX)_x$ perpendicular to $\omega$, and $V_x$ a general $i+1$-dimensional vector subspace of $(T^* X)_x$ containing $\omega$, the intersection of $\mathbb P (  (TY)_x) \times P ( (T^* X)_x)$ with a fixed $\dim X-2$-dimensional subspace of $\mathbb P ( (T^*X)_x)) \times \mathbb P( (TX)_x)$ is empty. ( The intersection of the strict transform of $\mathbb P(C)$ with the exceptional divisor has dimension one less than $\mathbb P(C)$, which itself has dimension one less than $C$, which has dimension $\dim X$.) To do this, it suffices to check that for each point $(v_1,v_2)$ of $\mathbb P ( (T^*X)_x) \times \mathbb P( (TX)_x)$ , the codimension of triples $\omega, (TY)_x $ such that $(v_1,v_2) \in \mathbb P (  (TY)_x) \times P ( (T^* X)_x)$ is at least $\dim X-1$.

%To do this, we can change the order to first choose $V_x$ generically, then $\omega$ generically in $V_x$, then $(TY)_x$ perpendicular to $\omega$. Note first that  that $V_x$ containing $v_2$ is a codimension $\dim X - (i+1)$ condition. Then if $v_1 \cdot \omega\neq 0$, then, $(TY)_x$ containing $v_1$ is a codimension $i$ condition, for a total of $\dim X-1$.  Alternatively, if $v_1 \cdot \omega=0$ is zero, then $(TY)_x$ containing $v_1$ is codimension $i-1$, but $v_1 \cdot \omega=0$ is a codimension $1$ condition unless the dot product with $v_1$ vanishes uniformly on $(TY)_x$, and that has codimension $1$ unless $i=0$ and $(v_1 \cdot v_2)=0$, but that is impossible as we assumed $i>0$. So we can save $1$ codimension this way but always lose one codimension in return, so the total codimension is $\dim X-1$, and thus the intersection is generically zero. 
 
 Finally, for $i = \dim X$, observe that any conical cycle whose projection to the cotangent space at $\inhyper{X}$ is the zero section was already the zero section, and any cycle whose restriction to a general hypersurface is the zero section was already the zero section.

\end{proof}



Recall the statement of Theorem \ref{first-Massey-intro}:

\begin{theorem}
\source{bounds-stalks-perverse-sheaves-characteristic-p-shende}
\notready\label{first-Massey}Let $X$ be a smooth variety over a perfect field $k$ and let $\ell$ be a prime invertible in $k$. Let $K$ be a perverse sheaf of $\mathbb F_\ell$-modules on $X$. 
\source{bounds-stalks-perverse-sheaves-characteristic-p-shende}

Then $\dim_{\mathbb F_\ell} \mathcal H^{-i}(K)_x$ is at most $i$th polar multiplicity of $CC(K)$ at $x$.  \end{theorem}



\begin{proof} This is an \'{e}tale-local question, so we may assume that $X$ is a smooth projective variety by passing to an affine open subset and embedding into projective space, then extending $K$ to keep it perverse. In fact, we fix an embedding into projective space. 

This follows by induction on $i$. Let $p$ and $q$ be the map defined by a general pencil of conics, as in Diagram \ref{pq-diagram}. We have a distinguished triangle \[ p^* K \to  R \Psi_q p^* K \to R \Phi_q p^*  K  .\] Taking stalk cohomology at $x$, we have an exact sequence

\[ \mathcal H^{-i-1}( R \Phi_q p^*  K)_x \to \mathcal H^{-i} ( K)_x \to \mathcal H^{-i} ( R \Psi_q p^* K)_x \to  \mathcal H^{-i}( R \Phi_q p^*  K)_x .\]


Because $K$ is perverse, $p^*K$ is perverse in a neighborhood of $x$, and thus $R \Psi_q p^* K[-1]$ is perverse.

Thus for $i=0$, $ \mathcal H^0 ( R \Psi_q p^* K)_x $ vanishes and we have
\[ \dim \mathcal H^0(K)_x \leq \dim  (R^{-1}\Phi_q p^*K)_x \leq \operatorname{dimtot} (R^{-1}  \Phi_q p^* K)_x = -  \operatorname{dimtot} (R \Phi_q p^* K)_x \] which is at most the multiplicity of the cotangent space at $x$ in $CC(K)$ which by Lemma \ref{multiplicity-hypersurface} is the $0$th polar multiplicity of $CC(K)$ in $x$.


By Lemma \ref{final-nearby-formula}(4), $\mathcal H^{i-1}( R \Phi_q p^*  K)_x$ vanishes unless $i=0$, so the map \[ \mathcal H^i ( K)_x \to \mathcal H^i ( R \Psi_q p^* K)_x \] is injective unless $i=0$. Thus for $i>0 $, we have

\[ \dim \mathcal H^{-i} ( (p^*K)_x)  \leq \mathcal H^{-i} ( R \Psi_q p^* K)_x. \]

Because $R \Psi_q p^* K [-1]$ is perverse, we can apply the induction hypothesis, to see that $\dim  \mathcal H^{-i} ( R \Psi_q p^* K)_x$  is at most the $i-1$st polar multiplicity of $CC( R \Psi_q p^* K [-1])$ at $x$. By Lemma \ref{final-nearby-formula}(1), $CC( R \Psi_q p^* K [-1])$ is the projection to the cotangent space of the conic of the restriction to the conic of $CC'(K)$. By Lemma \ref{multiplicity-hypersurface}, the $i-1$st polar multiplicity of this is the $i$th polar multiplicity of $CC(K)$, verifying the induction step.\end{proof}

In characteristic zero, the inequality $(R\Phi_f K)_x \leq \operatorname{dimtot} (R \Phi_f K)_x$ would be an identity, and we could use the Morse inequalities to derive additional information about the Betti numbers of $K$, as Massey does in \cite[Corollary 5.5]{Massey}, but in our case the analogue of the Morse inequalities are unhelpful.


\begin{proof}[Proof of Corollary \ref{vanishing-intro}] In view of Theorem \ref{first-Massey} it suffices to check that for $i< \dim X - \dim (SS(K)_x)$, the $i$th polar multiplicity of $CC(K)$ at $x$ vanishes.  For $V$ a vector bundle of rank $i+1$, $\mathbb P(V)$ has dimension $i$ in the fiber of $0$, and $\mathbb P( CC(K))$, which is contained in $\mathbb P( SS(K))$, has codimension $\dim X - \dim (SS(K)_x)$ in the fiber at zero, so for a generic $V$ these do not intersect and their intersection number, which is the polar multiplicity, vanishes. \end{proof}
